Questo lavoro di tesi ha affrontato il problema del rigging, skinning e retargeting automatico di animazioni su oggetti 3D inanimati, con particolare attenzione all'adattamento di animazioni antropomorfe a geometrie non convenzionali. L'obiettivo principale è stato quello di analizzare e confrontare diversi approcci automatici, valutandone l'efficacia all'interno di una pipeline basata su Blender.

I risultati ottenuti mostrano come le scelte effettuata nelle fasi di rigging e skinning influenzino in modo significativo la qualità e la controllabilità del retargeting delle animazioni. In particolare, l'approccio basato su UniRig si è dimostrato efficace nel mantenere una corrispondenza strutturale e semantica con le animazioni sorgente, facilitando il riutilizzo di dataset standard come Mixamo anche su oggetti non antropomorfi.

Al contrario, approcci più flessibili come RigAnything tendono a privilegiare la qualità visiva del movimento adattato alla geometria dell'oggetto, a scapito di una maggiore difficoltà nel controllo fine del retargeting e nella riproducibilità dei risultati su animazioni differenti.

Il presente lavoro presenta alcuni limiti. In primo luogo, la valutazione dei risultati è stata prevalentamente qualitativa, concentrandosi sull'analisi visiva e sulla coerenza del movimento, senza introdurre metriche quantitative standardizzate. Inoltre, i test sono stati condotti su un numero limitato di modelli e animazioni, scelti come casi di studio rappresentativi ma non esautivi.

Un ulteriore limite riguarda l'integrazione parziale di alcuni approcci automatici, che non ha consentito una valutazione uniforme di tutte le fasi della pipeline per ciascun metodo considerato.

Possibili sviluppi futuri di questo lavoro includono l'introduzione di metriche quantitative per la valutazione del retargeting, come l'analisi delle traiettorie dei punti articolari o la misurazione della deviazione rispetto a movimenti di riferimento. Tali metriche permetterebbero di affiancare all'analisi qualitativa una vautazione più oggettiva dei risultati.

Un ulteriore sviluppo potrebbe consistere nell'estensione dei test a un numero maggiori di modelli 3D e categorie di oggetti, al fine di valutare la robustezza degli approcci considerati in scenari più eterogenei

Infine, l'integrazione di tecniche di apprendimento automatico per l'ottimizzazione automatica del rigging e del retargeting potrebbe rappresentare una direzione promettente per migliorare ulteriormente l'adattamento delle animazioni a geometrie non convenzionali.